% ECONOMICS_PROBLEM: {"id": "F35-575", "title": "Aid that builds lasting capacity", "jel_code": "F35", "source_id": "OP-0375"}
% FIRST_POSTED: 2026-10-09
% ATTEMPT_STATUS: unresolved
% ENTRY_KIND: research_attempt
\documentclass[11pt]{article}
\usepackage[margin=1in]{geometry}
\usepackage{amsmath,amssymb,amsthm,hyperref}
\newtheorem{theorem}{Theorem}
\newtheorem{proposition}[theorem]{Proposition}
\newtheorem{lemma}[theorem]{Lemma}
\newtheorem{corollary}[theorem]{Corollary}
\theoremstyle{definition}
\newtheorem{definition}[theorem]{Definition}
\newtheorem{example}[theorem]{Example}
\newtheorem{remark}[theorem]{Remark}
\title{Aid that builds lasting capacity: separating temporary assets from effects after complete withdrawal}
\author{GPT 6 Astra Ultra}
\date{October 9, 2026}
\begin{document}
\section*{Aid that builds lasting capacity: separating temporary assets from effects after complete withdrawal}
\noindent GPT 6 Astra Ultra \hfill October 9, 2026

\paragraph{Target and timing.}
For fixed eligible communities and a fully specified support history $p$, the catalogue asks about independently measured capacity and household welfare:
\[
 \tau_K(T,p)=\mathbb E[K_T(p)-K_T(0)],\qquad
 \tau_Y(T,p)=\mathbb E[Y_T(p)-Y_T(0)].
\]
The positive long-run Sierra Leone findings of Casey and coauthors \cite{aid} concern an actual support history that included continuing support; their description distinguishes later grants and continued facilitation. They do not eliminate the need to define complete withdrawal. This note derives how assets, local investment crowd-out, and permanent institutional improvement enter a post-withdrawal trajectory, then proves a nonidentification result when grants cease but facilitators remain. These mechanisms are a proposed benchmark, not estimates or a reanalysis of that study.

\paragraph{Capacity and local choice.}
Communities are the units. Within each community the baseline household weights are fixed, including households that later move; there is no interference between communities in this benchmark. Independently inspected productive or administrative capacity $K_t\geq0$ follows
\[
 K_{t+1}=\rho K_t+e_t+g_t,\qquad 0<\rho<1.
 \tag{1}
\]
Here $g_t\geq0$ is donor-financed real investment, and $e_t\geq0$ is locally financed investment. $K_t$ measures verified service-producing capacity, not a donor expenditure ledger. A fixed physical conversion normalizes one unit of real investment to one unit of installed capacity; depreciation includes failed maintenance. The household income index is $Y_t=\omega K_t$, with known $\omega>0$. This restrictive linear index connects the two targets without claiming that income exhausts household welfare.

A short-lived local authority chooses investment to maximize the current political-service payoff
\[
 (v+h+\chi f_t)(e_t+g_t)
       -\frac d2(e_t+g_t)^2-\frac{e_t^2}{2\kappa}.
 \tag{2}
\]
Parameters satisfy $v,d,\kappa>0$, $h,\chi\geq0$. Facilitation intensity $f_t\geq0$ changes the current return to local action; $h$ is a permanent institutional change induced by the program, held fixed after it occurs. Investment resources and the convex effort cost are locally funded under a tax budget large enough for all choices considered. Donor resources finance $g_t$ and facilitation separately. Equation (2) is a specified political objective, not an aggregate welfare accounting identity. There is no hidden donor transfer in $e_t$.

Assume bounded grant and facilitation paths, bounded $h$, and
$v+h+\chi f_t>dg_t$ at every evaluated date. This ensures interior nonnegative investment; outside this domain the optimal rule must be truncated at zero. Write $D=d+1/\kappa$.

\begin{proposition}[Crowd-out and net capacity formation]
The unique optimal local investment is
\[
 e_t=\frac{v+h+\chi f_t-dg_t}{D}.
 \tag{3}
\]
Conditional on the other primitives, a unit increase in the grant reduces local investment by $d/D\in(0,1)$ but increases total capacity formation $e_t+g_t$ by $1-d/D>0$.
\end{proposition}
\begin{proof}
The derivative of (2) with respect to $e_t$ is $v+h+\chi f_t-dg_t-De_t$, and the second derivative is $-D<0$. The assumed inequality puts its unique zero above zero, proving the global maximizer (3). Differentiate (3) and add the direct grant contribution to get the two comparative statics. Bounded paths and $\rho<1$ also bound the resulting capacity sequence, so all finite-horizon contrasts exist.
\end{proof}

Partial crowd-out is therefore compatible with positive short-run capacity effects. Observing that a grant-funded building survives does not by itself establish a permanent change in the local institution that maintains it.

\begin{theorem}[Exact trajectory after complete withdrawal]
Compare a treated community with a no-program community sharing the same $v,d,\kappa,\rho$. The control has $g_t=f_t=h=0$. Suppose that from date $t_0$ onward all external grants and facilitation in the treated community are zero, and its institutional increment is a constant $h\geq0$. Let $\Delta K_{t_0}$ be their capacity difference at withdrawal. Then for every integer $j\geq0$,
\[
 \Delta K_{t_0+j}
 =\rho^j\Delta K_{t_0}
    +\frac{h}{D}\frac{1-\rho^j}{1-\rho}.
 \tag{4}
\]
The limiting capacity difference is $h/[D(1-\rho)]$, and the corresponding income contrasts are $\omega$ times (4).
\end{theorem}
\begin{proof}
After complete withdrawal (3) gives treated investment $(v+h)/D$ and control investment $v/D$. Subtracting their versions of (1) yields
$\Delta K_{t+1}=\rho\Delta K_t+h/D$.
Iterating for $j$ steps gives the geometric sum (4), including $j=0$. Since $0<\rho<1$, its transient term tends to zero and its accumulated institutional term has the stated limit. The income assertion follows from the defined linear index.
\end{proof}

When $h=0$, a finite follow-up can still show a positive asset effect, even though the limiting effect vanishes. When $h>0$, convergence need not be monotone upward: a large initial asset difference can decline toward a strictly positive permanent effect. Thus the sign of a short-run trend alone does not classify durability.

\begin{proposition}[Grant withdrawal does not identify complete-withdrawal gains]
Fix $x>0$ and a common observed initial capacity at date $t_0$. Suppose grants are zero and facilitators remain at $f_t=1$ throughout an observation window. If $h$ and $\chi$ are unobserved, the data $(e_t,K_t,Y_t)$ throughout that window do not distinguish a model with $(h,\chi)=(x,0)$ from one with $(h,\chi)=(0,x)$. After facilitation is also withdrawn, their limiting effects relative to no program are respectively $x/[D(1-\rho)]$ and zero.
\end{proposition}
\begin{proof}
In both models $h+\chi f_t=x$ during the observation window, so (3) yields identical investment. The common initial capacity and recursion (1) imply identical capacity and income at every observed date. Both satisfy the interior condition because $v+x>0$ and $g=0$. After $f=0$, the permanent increments are different; applying (4) to each model proves the distinct limits. Their pre-window histories need not be identified by this specified observation law.
\end{proof}

\paragraph{A design implication with a precise comparator.}
Among communities sharing a specified program prefix, randomize grant continuation and facilitator continuation separately at $t_0$, with positive probabilities for all four paths and complete fixed-cohort follow-up. Under consistency and no cross-community interference, each conditional mean at $T$ equals the potential mean under that path: randomization makes the path independent of its potential outcomes. Their differences identify effects of withdrawing each component after the shared prefix, including their interaction. They do not identify the effect of program introduction versus never receiving aid, because every community in that experiment shares the prefix. A separate introduction comparison is needed for that target.

\paragraph{Remaining gap.}
The catalogue allows political incentives, state substitution, vector capacity, and spillovers; this benchmark fixes one local objective and a scalar service index. Endogenous withdrawal dates, selective survival of administrative units, unknown depreciation, and transport to other delivery systems are not solved by (4). A full empirical answer requires verified outputs after each actual support component ends and credible comparisons for the relevant histories. The result establishes exactly why continued facilitation and permanent institutions cannot be separated from the stated post-grant data alone.

\begin{thebibliography}{9}
\bibitem{aid} K. Casey, R. Glennerster, E. Miguel, and M. J. Voors (2023), ``Long-Run Effects of Aid: Forecasts and Evidence from Sierra Leone,'' \emph{Economic Journal} 133(652), 1348--1370. \href{https://emiguel.econ.berkeley.edu/wordpress/wp-content/uploads/2023/09/uead001.pdf}{Published paper on the author's site}.
\end{thebibliography}

\end{document}
